%! Graphing Quadratic Functions & Transformations — Unit 2, Lesson 2.1 — Slides (Beamer)
% ELEVEN frames, in the profile's order (LESSON_SHAPE.md §1): title -> learning targets (+ how
% today runs, 5/35/10/10) -> warm-up (ending in "hold on to this") -> hook (the vote, LEFT
% UNRESOLVED) -> four notes frames (rows 1-2 = section 1, rows 3-4 = section 2; the crux frame
% flagged \sectionlabel[redacc]{} with a block giving the case where the two answers disagree)
% -> Guided Practice -> debrief -> close & start the homework.
% There is no group activity and no individual-practice launch frame: the homework IS the
% individual practice, started in class. Beamer has no \CourseName — "Algebra 2" is literal.
% The header macro is \forestheader{Title}; \navyheader does not exist in this course.
\documentclass[aspectratio=169, 11pt]{beamer}
\usepackage{algebra2-beamer}   % provides \forestheader{} and \sectionlabel[color]{}
% Coordinate grid: {xmin}{xmax}{ymin}{ymax}
\newcommand{\qgrid}[4]{%
  \draw[step=1, linegray!40, very thin] (#1,#3) grid (#2,#4);
  \draw[->, thick, charcoal] (#1,0)--(#2,0) node[below right, font=\tiny]{$x$};
  \draw[->, thick, charcoal] (0,#3)--(0,#4) node[left, font=\tiny]{$y$};}
\newcommand{\fdot}[3]{\fill[#1] (#2,#3) circle (3.2pt);}

\begin{document}

% TITLE SLIDE (CourseName is not defined in beamer, so the name is literal)
{
\setbeamercolor{background canvas}{bg=forest}
\begin{frame}[plain]
  \vspace{2.8cm}\hspace{0.55cm}%
  \begin{minipage}{0.9\textwidth}
    {\color{goldacc}\bfseries\small LESSON 2.1}\\[0.5cm]
    {\color{white}\bfseries\LARGE Graphing Quadratic Functions \& Transformations}\\[1.2cm]
    {\color{forestmid}\small Algebra 2~$\cdot$~Unit 2: Quadratic Functions}
  \end{minipage}
\end{frame}
}

% LEARNING TARGETS
\begin{frame}[t, plain]
  \forestheader{Today's targets}
  \sectionlabel{What you will be able to do}
  \begin{itemize}\setlength{\itemsep}{6pt}
    \item Describe any parabola as a \textbf{transformation} of the \textbf{parent function}
          $y=x^2$: \textbf{shift}, \textbf{stretch}, \textbf{reflection}.
    \item Read the \textbf{vertex} $(h,k)$ and the \textbf{axis of symmetry} $x=h$ straight out
          of \textbf{vertex form} $a(x-h)^2+k$.
    \item Read the $y$-intercept out of \textbf{standard form} $ax^2+bx+c$ and its axis from
          $x=-\frac{b}{2a}$.
    \item Read the \textbf{zeros} out of \textbf{intercept (factored) form} $a(x-p)(x-q)$ and
          put the axis at their midpoint.
  \end{itemize}
  \vspace{4pt}
  \begin{block}{How today runs}
    Warm-up (5) $\to$ notes and one curve together (35) $\to$ debrief (10) $\to$
    \textbf{start the homework, on your own} (10).
  \end{block}
\end{frame}

% WARM-UP
\begin{frame}[t, plain]
  \forestheader{Warm-Up: three things you already know}
  \sectionlabel{5 minutes --- alone, then check with a neighbour}
  \begin{enumerate}\setlength{\itemsep}{7pt}
    \item \textbf{Multiply it out.} Expand $(x-1)^2-4$, then expand $(x+1)(x-3)$. Did they give
          the same answer?
    \item \textbf{Moving a parent graph (Lesson 1.3).} $y=|x|+2$ moves the V which way?
          $y=|x-4|$? $y=-|x|$?
    \item \textbf{Two parabolas.} The parent $y=x^2$, and $y=(x-3)^2$ beside it.
  \end{enumerate}
  \vspace{6pt}
  \begin{block}{Hold on to this}
    $y=(x-3)^2$ moved the parent \emph{right} $3$ --- even though the sign inside says
    ``minus.'' Keep that question: why does subtracting inside move it the \emph{other} way?
  \end{block}
\end{frame}

% HOOK — left unresolved
\begin{frame}[t, plain]
  \forestheader{Hook: one curve in three costumes}
  \sectionlabel[goldacc]{The same parabola you read in Lesson 2.0}
  \begin{columns}[T]
    \begin{column}{0.54\textwidth}
      \[ x^2-2x-3 \]
      \[ (x-1)^2-4 \]
      \[ (x+1)(x-3) \]
      Three expressions --- and \emph{one} graph on the right.
    \end{column}
    \begin{column}{0.42\textwidth}\centering
      \begin{tikzpicture}[scale=0.34]
        \qgrid{-3.5}{5.5}{-5.5}{5.5}
        \begin{scope}
          \clip (-3.5,-5.5) rectangle (5.5,5.5);
          \draw[very thick, forest, domain=-2.4:4.4, samples=80, smooth, variable=\x]
            plot (\x,{(\x-1)^2-4});
        \end{scope}
        \fdot{forest}{1}{-4}
      \end{tikzpicture}
    \end{column}
  \end{columns}
  \begin{block}{Vote --- and we do not settle it yet}
    \textbf{Circle one: three different parabolas \; / \; one parabola wearing three costumes.}
    \quad And: which costume tells you the \textbf{vertex} fastest? Which one the \textbf{zeros}?
  \end{block}
\end{frame}

% NOTES ROW 1 — THE PARENT AND ITS SHIFTS  (section 1a)
\begin{frame}[t, plain]
  \forestheader{Notes: the parent $y=x^2$ and its shifts}
  \sectionlabel{notes, Its Shifts $\cdot$ I do problem 1, you do problem 2}
  \begin{columns}[T]
    \begin{column}{0.50\textwidth}
      \begin{itemize}\setlength{\itemsep}{4pt}
        \item $f(x)+k$: a number \emph{outside} slides it \textbf{up or down}. $y=x^2+3$ ---
              vertex $(0,3)$.
        \item $f(x-h)$: a change \emph{inside} slides it \textbf{left or right}. $y=(x-2)^2$ ---
              vertex $(2,0)$.
        \item So $(x-2)^2$ moves the parent \textbf{right} $2$.
      \end{itemize}
    \end{column}
    \begin{column}{0.46\textwidth}\centering
      \begin{tikzpicture}[scale=0.32]
        \qgrid{-3.5}{5.5}{-1.5}{8.5}
        \begin{scope}
          \clip (-3.5,-1.5) rectangle (5.5,8.5);
          \draw[very thick, charcoal, domain=-3.2:3.2, samples=80, smooth, variable=\x]
            plot (\x,{(\x)^2});
          \draw[very thick, forest, domain=-2.6:2.6, samples=80, smooth, variable=\x]
            plot (\x,{(\x)^2+3});
          \draw[very thick, redacc, domain=-1.2:5.2, samples=80, smooth, variable=\x]
            plot (\x,{(\x-2)^2});
        \end{scope}
        \fdot{charcoal}{0}{0} \fdot{forest}{0}{3} \fdot{redacc}{2}{0}
      \end{tikzpicture}
    \end{column}
  \end{columns}
  \begin{block}{Subtracting inside moves it the other way}
    $(x-2)^2$ is zero when $x=2$, not when $x=-2$. The bracket asks what you must subtract to
    get back to the parent's input --- so the graph slides the \emph{opposite} way to the sign.
  \end{block}
\end{frame}

% NOTES ROW 2 — STRETCH AND FLIP  (section 1b)
\begin{frame}[t, plain]
  \forestheader{Notes: stretch, squash, flip --- that is $a$}
  \sectionlabel{notes, Stretch and Flip $\cdot$ I do problem 3, you do 4--5}
  \begin{columns}[T]
    \begin{column}{0.48\textwidth}
      \begin{itemize}\setlength{\itemsep}{4pt}
        \item $|a|>1$: \textbf{narrower} than the parent.
        \item $0<|a|<1$: \textbf{wider}.
        \item $a<0$: \textbf{reflected} --- opens down, so the vertex is a \textbf{maximum}.
      \end{itemize}
      \vspace{3pt}
      Check it with numbers, not with your eyes: at $x=2$, $y=x^2$ gives $4$ and $y=3x^2$ gives
      $12$.
    \end{column}
    \begin{column}{0.48\textwidth}\centering
      \begin{tikzpicture}[scale=0.32]
        \qgrid{-4.5}{4.5}{-4.5}{8.5}
        \begin{scope}
          \clip (-4.5,-4.5) rectangle (4.5,8.5);
          \draw[very thick, charcoal, domain=-3.2:3.2, samples=80, smooth, variable=\x]
            plot (\x,{(\x)^2});
          \draw[very thick, forest, domain=-1.9:1.9, samples=80, smooth, variable=\x]
            plot (\x,{3*(\x)^2});
          \draw[very thick, redacc, domain=-3.2:3.2, samples=80, smooth, variable=\x]
            plot (\x,{-0.5*(\x)^2});
        \end{scope}
        \fdot{charcoal}{0}{0}
      \end{tikzpicture}
    \end{column}
  \end{columns}
  \begin{block}{Same vertex, three widths}
    $y=x^2$, $y=3x^2$, $y=-\tfrac12 x^2$ all have vertex $(0,0)$. Only $a$ changed --- and the
    red one opens \emph{down}, so $(0,0)$ went from a minimum to a maximum.
  \end{block}
\end{frame}

% NOTES ROW 3 — VERTEX FORM READS ITSELF  (section 2a)
\begin{frame}[t, plain]
  \forestheader{Notes: vertex form reads itself}
  \sectionlabel{notes, Vertex Form $\cdot$ I do problem 6, you do 7--8}
  \[ y=a(x-h)^2+k \quad\Longrightarrow\quad \text{vertex } (h,k), \quad \text{axis } x=h \]
  \begin{itemize}\setlength{\itemsep}{4pt}
    \item Read $a$ (direction and width) $\to$ read $h$, \textbf{flipping the sign inside} $\to$
          read $k$, no flip.
    \item Our anchor: $y=(x-1)^2-4$. So $a=1$, $h=1$, $k=-4$: vertex $(1,-4)$, axis $x=1$,
          opens up.
    \item Plot the vertex, then use \textbf{symmetry} --- a point one step left has a twin the
          same height one step right.
  \end{itemize}
  \begin{block}{No work required}
    Vertex form is the only one of the three that hands you the vertex with no arithmetic at all.
    That is what makes it worth recognising on sight.
  \end{block}
\end{frame}

% NOTES ROW 4 — THE CRUX  (section 2b)
\begin{frame}[t, plain]
  \forestheader{Notes: three costumes, one curve --- and the bracket that lies}
  \sectionlabel[redacc]{notes, One Curve $\cdot$ the crux, problems 9--10 --- which form you are in decides}
  \begin{columns}[T]
    \begin{column}{0.48\textwidth}
      \textbf{Standard} $x^2-2x-3$: $y$-intercept $(0,-3)$ free; axis
      $x=-\frac{-2}{2(1)}=1$; then $y=-4$. \\[4pt]
      \textbf{Intercept} $(x+1)(x-3)$: zeros $-1$ and $3$; axis $\frac{-1+3}{2}=1$. \\[4pt]
      \textbf{Vertex} $(x-1)^2-4$: vertex $(1,-4)$, axis $x=1$. \\[4pt]
      \textbf{All three agree --- it is one curve.}
    \end{column}
    \begin{column}{0.48\textwidth}
      \begin{block}{Where the two answers disagree}
        In \emph{vertex} form, $(x+4)^2$ shifts the parent \textbf{left} $4$. \\[3pt]
        As a \emph{factor}, $(x+4)$ gives the zero $x=\mathbf{-4}$ --- and in $(x+1)(x-3)$ the
        zeros are $-1$ and $3$, not $1$ and $-3$. \\[3pt]
        \textbf{Same bracket, opposite-looking jobs. Which form you are in decides.}
      \end{block}
    \end{column}
  \end{columns}
\end{frame}

% GUIDED PRACTICE — WE DO
\begin{frame}[t, plain]
  \forestheader{Guided Practice: we do this one together}
  \sectionlabel[goldacc]{You hold the pen --- problems 11--14: one new curve, three costumes}
  $y=-(x-2)^2+9 \;=\; -x^2+4x+5 \;=\; -(x+1)(x-5)$ --- the graph is on your page.
  \begin{enumerate}\setlength{\itemsep}{4pt}
    \item[\textbf{11.}] Vertex form: vertex, axis, up or down --- and is $9$ a maximum or a
          minimum?
    \item[\textbf{12.}] Standard form: the $y$-intercept straight off, then the axis from
          $x=-\frac{b}{2a}$.
    \item[\textbf{13.}] Intercept form: the zeros, and the axis at their midpoint.
    \item[\textbf{14.}] You found the axis three ways. Do they agree, and why must they? A
          classmate read the zeros as $1$ and $-5$ --- what did they do wrong?
  \end{enumerate}
  \vspace{3pt}
  \begin{block}{The questions I will ask}
    ``Which form is this, and what does it hand you for free?'' \quad ``What is $h$ --- and did
    you flip the sign?'' \quad ``$(x+1)$ is zero at which $x$?''
  \end{block}
\end{frame}

% DEBRIEF
\begin{frame}[t, plain]
  \forestheader{Debrief: pens down, papers up}
  \sectionlabel[redacc]{10 minutes --- aloud, nothing to fill in}
  \begin{enumerate}\setlength{\itemsep}{5pt}
    \item \textbf{One curve, three costumes.} The vote from the hook is settled: $x^2-2x-3$,
          $(x-1)^2-4$ and $(x+1)(x-3)$ are the \emph{same} parabola. Each form just hands you a
          different feature first.
    \item \textbf{The sign flip.} $(x-1)^2$ moves the parent \emph{right} $1$; $x^2-1$ moves it
          \emph{down} $1$. Inside versus outside, every time.
    \item \textbf{The same bracket, two forms.} $(x+4)^2$ shifts \emph{left} $4$; the factor
          $(x+4)$ gives the zero $-4$. Which form you are in decides what the bracket means.
    \item \textbf{$a$ does two jobs.} Direction (sign) and width (size) --- and when $a<0$ the
          vertex is a \emph{maximum}, which is Lesson~2.0's language read off the equation.
  \end{enumerate}
  \vspace{3pt}
  \begin{block}{Say it without the notes}
    ``$y=(x-1)^2-4$: vertex $(1,-4)$, axis $x=1$, opens up --- and I flipped the sign inside to
    get $h=1$.''
  \end{block}
\end{frame}

% CLOSE & START HOMEWORK
\begin{frame}[t, plain]
  \forestheader{Close --- and start the homework, on your own}
  \sectionlabel{10 minutes $\cdot$ silent $\cdot$ the Homework page in your packet}
  You walked in able to \emph{read} a parabola someone drew for you. You walk out able to take an
  equation, name which of the three forms it is in, and say where its vertex, axis, zeros and
  $y$-intercept are --- and how the parent $y=x^2$ was moved to get there.
  \vspace{5pt}
  \begin{block}{Homework --- scored, due the first class after two study halls}
    We do item~1(a) together right now; then you keep going, alone. Six items: reading vertex
    form, the inside-versus-outside pair, a graph to read \emph{and} write as a rule, intercept
    form and its special case, a stone arch, and one multiple-choice item. Whatever is left is
    due the first class after two study halls --- I will announce the date in class and on
    TurtleNet.
  \end{block}
  \vspace{3pt}
  \textbf{Next:} Lesson 2.2 --- where the factored form \emph{comes from}, and how to build it
  yourself.
\end{frame}

\end{document}
